% Hao Wang — CV (English)
% Black-and-white print layout; typography matches https://wanghao9610.github.io/ :
% serif headings (Source Serif Pro) with sans body (Source Sans Pro).
% Build: make   (latexmk -xelatex)

\documentclass[10pt,letterpaper]{article}

\usepackage[margin=0.62in,top=0.55in,bottom=0.55in]{geometry}
\usepackage{fontspec}
\setmainfont{SourceSerifPro-Regular.otf}[
  BoldFont=SourceSerifPro-Semibold.otf,
  ItalicFont=SourceSerifPro-RegularIt.otf,
  BoldItalicFont=SourceSerifPro-SemiboldIt.otf]
\setsansfont{SourceSansPro-Regular.otf}[
  BoldFont=SourceSansPro-Semibold.otf,
  ItalicFont=SourceSansPro-RegularIt.otf,
  BoldItalicFont=SourceSansPro-SemiboldIt.otf]
\renewcommand{\familydefault}{\sfdefault}
\usepackage{xeCJK}
\setCJKmainfont{FandolSong-Regular.otf}[BoldFont=FandolSong-Bold.otf]
\usepackage{xcolor}
\usepackage{hyperref}
\usepackage{enumitem}
\usepackage{tabularx}
\usepackage{titlesec}
\usepackage{array}
\usepackage{ragged2e}
\usepackage{microtype}
\usepackage{needspace}
\usepackage{fontawesome5}
\usepackage{etoolbox}

% ---- monochrome print palette --------------------------------------------
\definecolor{ink}{HTML}{000000}
\definecolor{text}{HTML}{000000}
\definecolor{muted}{HTML}{3A3A3A}
\definecolor{soft}{HTML}{555555}
\definecolor{accent}{HTML}{000000}
\definecolor{rule}{HTML}{BFBFBF}

\hypersetup{
  colorlinks=true,
  urlcolor=accent,
  linkcolor=accent,
  pdftitle={Hao Wang - CV},
  pdfauthor={Hao Wang}
}

\pagestyle{empty}
\color{text}
\setlength{\parindent}{0pt}
\setlength{\parskip}{0pt}
\setlist[itemize]{leftmargin=1.15em,itemsep=1.6pt,topsep=2pt,parsep=0pt,label=\textcolor{accent}{\scriptsize\raisebox{.25ex}{$\bullet$}}}
\renewcommand{\arraystretch}{1.06}
\setlength{\emergencystretch}{1.5em}

% ---- section: serif title + hairline ---------------------------------------
\titleformat{\section}
  {\rmfamily\Large\bfseries\color{ink}}
  {}
  {0pt}
  {}
  [{\vspace{-3pt}\color{rule}\titlerule[0.6pt]}]
\titlespacing*{\section}{0pt}{9pt}{5pt}

\newcommand{\datebadge}[1]{\textcolor{accent}{\footnotesize\bfseries #1}}
\newcommand{\metaline}[1]{\textcolor{muted}{\small #1}}
\newcommand{\venue}[1]{\textcolor{accent}{\footnotesize\bfseries #1}}
\newcommand{\stars}[1]{\ifstrempty{#1}{}{\hfill\textcolor{soft}{\footnotesize\faStar\,#1}}}
\newcommand{\sep}{{\color{rule}\;\textbar\;}}

% \cvitem{title}{date}{subtitle}{body}
\newcommand{\cvitem}[4]{%
  \Needspace{6\baselineskip}
  \vspace{3pt}
  \begin{tabularx}{\textwidth}{@{}>{\RaggedRight\arraybackslash}X>{\RaggedLeft\arraybackslash}p{0.25\textwidth}@{}}
    \textbf{\textcolor{ink}{#1}} & \datebadge{#2}
  \end{tabularx}
  \vspace{-1pt}
  \metaline{#3}\par
  \vspace{1pt}
  #4
}

% \publication{url}{title}{venue}{authors}{summary}
\newcommand{\publication}[5]{%
  \Needspace{5\baselineskip}
  \vspace{2pt}
  \begin{tabularx}{\textwidth}{@{}>{\RaggedRight\arraybackslash}X>{\RaggedLeft\arraybackslash}p{0.16\textwidth}@{}}
    \textbf{\href{#1}{\textcolor{ink}{#2}}} & \venue{#3}
  \end{tabularx}
  \vspace{-1pt}
  {\small #4}\par
  \vspace{1pt}
  \textcolor{muted}{\small #5}\par
  \vspace{4pt}
}

% \repo{url}{name}{tag}{stars}{description}
\newcommand{\repo}[5]{%
  \textbf{\href{#1}{\textcolor{ink}{#2}}}\ \venue{#3}\stars{#4}\par
  \vspace{-1pt}
  \textcolor{muted}{\small #5}\par
  \vspace{3.5pt}
}

\begin{document}

% ============================ Header =========================================
\begin{minipage}{\textwidth}
  \vspace{2pt}
  \centering
  {\rmfamily\fontsize{28}{32}\selectfont\bfseries\color{ink} Hao Wang
   {\fontsize{20}{24}\selectfont\mdseries\color{muted} 王豪}}\par
  \vspace{5pt}
  {\rmfamily\itshape\large\color{muted} Leave the world a little different.}\par
  \vspace{4pt}
  {\normalsize\color{accent} Ph.D. Candidate · Computer Vision · Multimodal AI · Agentic AI}\par
  \vspace{7pt}
  {\small\color{muted}
    \faEnvelope\ \href{mailto:wanghao9610@gmail.com}{wanghao9610@gmail.com}
    \sep
    \faGlobe\ \href{https://wanghao9610.github.io/}{wanghao9610.github.io}
    \sep
    \faGithub\ \href{https://github.com/wanghao9610}{wanghao9610}
    \sep
    \faWeixin\ wangh9610
    \sep
    \faMapMarker*\ Shenzhen, China
  }\par
  \vspace{7pt}
  {\color{ink}\rule{\linewidth}{0.8pt}}\par
  \vspace{5pt}
  \begin{minipage}[t]{0.62\textwidth}
    \RaggedRight\small\color{muted} HCP Lab, Sun Yat-sen University \& Pengcheng Laboratory
  \end{minipage}%
  \hfill
  \begin{minipage}[t]{0.36\textwidth}
    \RaggedLeft\small\color{muted} Expected graduation: \textbf{\color{accent}December 2026}
  \end{minipage}\par
  \vspace{1pt}
  {\small\color{muted}\RaggedRight Open to industry research roles and research collaborations.\par}
\end{minipage}

% ============================ Profile ========================================
\section{Profile}
Ph.D. candidate whose research centers on \textbf{open-ended computer vision} and \textbf{multimodal large language models}, and who is increasingly exploring \textbf{multimodal agentic models}. First-author work includes X2SAM (ECCV 2026), X-SAM (AAAI 2026), OV-DINO, and TMANet (ICIP 2021), spanning unified image/video segmentation, open-vocabulary detection, and temporal pixel-level understanding. All representative works are open-sourced with project pages and code (1.1k+ GitHub stars in total). Interested in connecting foundation models with reliable dense perception and agentic workflows for practical vision systems.

\section{Research Highlights}
\begin{itemize}
  \item Built a first-author research line on open-ended perception, from video semantic segmentation to open-vocabulary detection and multimodal any-segmentation.
  \item Developed multimodal segmentation MLLMs that accept both natural-language instructions and visual prompts, bridging high-level reasoning and dense mask prediction.
  \item Extended any-segmentation from images to videos with a mask-memory mechanism for temporally consistent object- and pixel-level understanding.
  \item Maintain open-source research code and an agent-oriented research toolchain (STAR / STORY / STAGE), emphasizing reproducible and open research practice.
\end{itemize}

% ============================ Publications ===================================
\section{First-Author Publications}
\publication
  {https://wanghao9610.github.io/X2SAM/}
  {X2SAM: Any Segmentation in Images and Videos}
  {ECCV 2026}
  {\textbf{Hao Wang}, Limeng Qiao, Chi Zhang, Guanglu Wan, Lin Ma, Xiangyuan Lan, Xiaodan Liang}
  {A unified segmentation MLLM that extends any-segmentation from images to videos, supporting conversational instructions and visual prompts through Mask Memory for temporally consistent pixel-level perception. \href{https://wanghao9610.github.io/X2SAM/}{Project} \sep \href{https://arxiv.org/abs/2605.00891}{Paper} \sep \href{https://github.com/wanghao9610/X2SAM}{Code}}

\publication
  {https://wanghao9610.github.io/X-SAM/}
  {X-SAM: From Segment Anything to Any Segmentation}
  {AAAI 2026}
  {\textbf{Hao Wang}, Limeng Qiao, Zequn Jie, Zhijian Huang, Chengjian Feng, Qingfang Zheng, Lin Ma, Xiangyuan Lan, Xiaodan Liang}
  {A unified MLLM framework that extends segment-anything capabilities to any segmentation, enabling instruction-driven interaction and pixel-level perceptual understanding. \href{https://wanghao9610.github.io/X-SAM/}{Project} \sep \href{https://arxiv.org/abs/2508.04655}{Paper} \sep \href{https://github.com/wanghao9610/X-SAM}{Code}}

\publication
  {https://wanghao9610.github.io/OV-DINO/}
  {OV-DINO: Unified Open-Vocabulary Detection with Language-Aware Selective Fusion}
  {arXiv 2024}
  {\textbf{Hao Wang}, Pengzhen Ren, Zequn Jie, Xiao Dong, Chengjian Feng, Yinlong Qian, Lin Ma, Dongmei Jiang, Yaowei Wang, Xiangyuan Lan, Xiaodan Liang}
  {A unified open-vocabulary detector pre-trained on diverse large-scale datasets with language-aware selective fusion, improving category generalization through visual-language alignment. \href{https://wanghao9610.github.io/OV-DINO/}{Project} \sep \href{https://arxiv.org/abs/2407.07844}{Paper} \sep \href{https://github.com/wanghao9610/OV-DINO}{Code}}

\publication
  {https://arxiv.org/abs/2102.08643}
  {TMANet: Temporal Memory Attention for Video Semantic Segmentation}
  {ICIP 2021}
  {\textbf{Hao Wang}, Weining Wang, Jing Liu}
  {A temporal memory attention method for long-range video semantic segmentation that captures cross-frame relations without optical-flow prediction. \href{https://arxiv.org/abs/2102.08643}{Paper} \sep \href{https://github.com/wanghao9610/TMANet}{Code}}

% ============================ Open source =====================================
\Needspace*{16\baselineskip}
\section{Open-Source Projects}
\metaline{Research code and tooling on \href{https://github.com/wanghao9610}{GitHub}; star counts as of September 2026.}\par
\vspace{4pt}
\begin{minipage}[t]{0.485\textwidth}
  \repo{https://github.com/wanghao9610/OV-DINO}{OV-DINO}{arXiv 2024}{414}{Unified open-vocabulary detection with language-aware selective fusion.}
  \repo{https://github.com/wanghao9610/X-SAM}{X-SAM}{AAAI 2026}{391}{From segment anything to any segmentation.}
  \repo{https://github.com/wanghao9610/TMANet}{TMANet}{ICIP 2021}{128}{Temporal memory attention for video semantic segmentation.}
  \repo{https://github.com/wanghao9610/X2SAM}{X2SAM}{ECCV 2026}{104}{Any segmentation in images and videos with one unified MLLM.}
\end{minipage}%
\hfill
\begin{minipage}[t]{0.485\textwidth}
  \repo{https://github.com/wanghao9610/STAR}{STAR}{Harness · WIP}{51}{Systematic Toolchain for AI Research: an agent harness for planning, running, and tracking research.}
  \repo{https://github.com/wanghao9610/STORY}{STORY \,/\, \href{https://github.com/wanghao9610/STAGE}{STAGE}}{Harness · WIP}{}{Companion toolchains for organizing research over years and for authoring, guiding, and editing.}
  \repo{https://omdsh-plugins.github.io/}{omdsh-plugins}{Collection}{}{Oh My DeepSeek Harness: a curated plugin collection for the DeepSeek Harness.}
  \repo{https://github.com/wanghao9610/arXivTeX}{arXivTeX}{Template}{33}{Paper template for arXiv or any conference. Also: \href{https://github.com/wanghao9610/Claude-Model-Proxy}{Claude Model Proxy}, \href{https://github.com/wanghao9610/Codex-Remote-Connector}{Codex Remote Connector}.}
\end{minipage}

% ============================ Experience =====================================
\section{Research Experience}
\cvitem
  {Meituan · M17-MM}
  {2025.01 -- Present}
  {Research Intern}
  {\begin{itemize}
    \item Led X2SAM, a unified any-segmentation MLLM for images and videos with conversational instructions and visual prompts.
    \item Designed a mask-memory mechanism that preserves object-level consistency in video segmentation and interactive perception; owned training, evaluation, and open-source release.
  \end{itemize}}

\cvitem
  {Meituan · Vision Intelligence Department}
  {2022.07 -- 2025.01}
  {Research Intern}
  {\begin{itemize}
    \item Led OV-DINO and X-SAM, advancing open-vocabulary recognition and any segmentation.
    \item Studied language-aware fusion and multimodal supervision for scalable visual recognition and segmentation.
  \end{itemize}}

\cvitem
  {Tencent · AI Platform Department}
  {2021.05 -- 2021.08}
  {Applied Research Intern}
  {}

\cvitem
  {Huawei · Photo Processing Department}
  {2019.09 -- 2020.07}
  {Applied Project Intern}
  {}

% ============================ Education ======================================
\section{Education}
\cvitem
  {Sun Yat-sen University \& Pengcheng Laboratory}
  {2022.09 -- Present}
  {Ph.D. student, School of Intelligent Systems Engineering}
  {\begin{itemize}
    \item Co-advised by Prof. Xiaodan Liang and Assoc. Prof. Xiangyuan Lan.
    \item Research: image/video segmentation, open-ended visual perception, multimodal foundation models, and multimodal agentic models. Expected graduation: December 2026.
  \end{itemize}}

\cvitem
  {University of Chinese Academy of Sciences \& Institute of Automation, CAS}
  {2019.09 -- 2022.06}
  {Master student, School of Artificial Intelligence; advised by Prof. Jing Liu}
  {}

\cvitem
  {Beijing Jiaotong University}
  {2015.09 -- 2019.06}
  {Bachelor student, School of Electronic and Information Engineering}
  {}

% ============================ Misc ===========================================
\section{Co-Authored Publication}
\publication
  {https://ieeexplore.ieee.org/abstract/document/10096063}
  {WL-MSR: Watch and Listen for Multimodal Subtitle Recognition}
  {ICASSP 2023}
  {Jiawei Liu, \textbf{Hao Wang}, Weining Wang, Xingjian He, Jing Liu}
  {Transformer-based multimodal subtitle recognition that fuses OCR and ASR information with mask/crop strategies and multi-level identity embeddings.}

\Needspace*{5\baselineskip}
\section{Service \& Awards}
\begin{tabular}{@{}>{\bfseries\color{ink}}l@{\hspace{0.08\textwidth}}l@{}}
Conference Reviewer & AAAI 2026, ECCV 2024, ICCV 2023 \\
Journal Reviewer & Proceedings of the IEEE \\
Award & 1st place, the 1st VSPW Challenge Workshop, ICCV 2021 (2021.09)
\end{tabular}

\end{document}
